\clearpage
\VerdierPhysicalPageAnchor{0209}%
\begingroup\small\setlength{\parskip}{0pt}\setlength{\emergencystretch}{1.5em}
\setlength{\abovedisplayskip}{4pt plus 2pt minus 2pt}%
\setlength{\belowdisplayskip}{4pt plus 2pt minus 2pt}%

\VerdierSemanticAnchor{ch3d-3-1-3-7}
\[
\textit{(7) $K^{b}(\mathcal A)$ の任意の $X$ と
$K^{+}(\mathcal A)$ の任意の $Y$ に対して、}\qquad
\operatorname{Ext}_{\mathcal A}^{0}(X,Y)\simeq
\varinjlim_{(\operatorname{Qis}^{b}/X)^o}
\operatorname{Hom}_{K(\mathcal A)}(\mathord\cdot,Y).
\]

\VerdierSemanticAnchor{ch3d-3-1-3-8}
\[
\textit{(8) $K^{-}(\mathcal A)$ の任意の $X$ と
$K^{b}(\mathcal A)$ の任意の $Y$ に対して、}\qquad
\operatorname{Ext}_{\mathcal A}^{0}(X,Y)\simeq
\varinjlim_{Y\backslash\operatorname{Qis}^{b}}
\operatorname{Hom}_{K(\mathcal A)}(X,\mathord\cdot).
\]

\VerdierSemanticAnchor{ch3d-3-1-3-proof}
式 (1) は局所化圏の定義（第 II 章，2.2.2）から従う。式 (2) は式 (1)
から反対圏へ移ることによって従う（第 II 章，2.2.4）。したがって式 (3)
は式 (1), (2) から従う。式 (4) は定義である（第 I 章，1.3.8）。式 (5),
(6) は命題 (1.2.1) の b), b)\textsuperscript{bis} から従い、式 (7),
(8) は同命題の c), c)\textsuperscript{bis} から従う。

\VerdierSemanticAnchor{ch3d-3-1-4}
\medskip\noindent\textbf{命題 3.1.4.}\quad
\textit{圏 $\mathcal A$ の各対象が、$\mathcal A$ のある入射対象の部分対象
（それぞれ、ある射影対象の商対象）と同型であるとする。}

\VerdierSemanticAnchor{ch3d-3-1-4-1}
\noindent\textit{1) $K^{+}(\mathcal A)$ の各対象 $Y$（それぞれ
$K^{-}(\mathcal A)$ の各対象 $X$）は、入射型の右分解
$u\colon Y\longrightarrow I^\bullet$（それぞれ射影型の左分解
$u\colon P^\bullet\longrightarrow X$）をもつ (1.1.4)。}

\VerdierSemanticAnchor{ch3d-3-1-4-2}
\noindent\textit{2) この分解は同型を定める：}
\[
\operatorname{Ext}_{\mathcal A}^{0}(X,Y)\simeq
\operatorname{Hom}_{K(\mathcal A)}(X,I^\bullet),
\qquad\textit{$K(\mathcal A)$ の任意の $X$ に対して}
\]
\noindent\textit{（それぞれ、}
\[
\operatorname{Ext}_{\mathcal A}^{0}(X,Y)\simeq
\operatorname{Hom}_{K(\mathcal A)}(P^\bullet,Y),
\qquad\textit{$K(\mathcal A)$ の任意の $Y$ に対して）。}
\]

\VerdierSemanticAnchor{ch3d-3-1-4-3}
\noindent\textit{3) $\mathcal A$ のコホモロジー次元は、$\mathcal A$ の
対象の入射次元
（それぞれ射影次元）の上限である。}

\VerdierSemanticAnchor{ch3d-3-1-4-4}
\noindent\textit{4) $\mathcal A$ のコホモロジー次元が有限であるとき、
$K(\mathcal A)$ の各対象 $Y$（それぞれ各対象 $X$）は、入射型の右分解
$u\colon Y\longrightarrow I^\bullet$（それぞれ射影型の左分解
$u\colon P^\bullet\longrightarrow X$）をもつ。}

\VerdierSemanticAnchor{ch3d-3-1-4-5}
\noindent\textit{5) 4) の仮定と記法のもとで、このような分解は同型を定める：}
\[
\operatorname{Ext}_{\mathcal A}^{0}(X,Y)\simeq
\operatorname{Hom}_{K(\mathcal A)}(X,I^\bullet),
\qquad\textit{$K(\mathcal A)$ の任意の $X$ に対して}
\]
\noindent\textit{（それぞれ、}
\[
\operatorname{Ext}_{\mathcal A}^{0}(X,Y)\simeq
\operatorname{Hom}_{K(\mathcal A)}(P^\bullet,Y),
\qquad\textit{$K(\mathcal A)$ の任意の $Y$ に対して）。}
\]
\endgroup
